% =====================================================================
% PLANTILLA: Evaluación -- Área de Matemáticas
% María Mercedes Cantillo · Colegio San Alberto Magno
% ---------------------------------------------------------------------
% 1. Copia este archivo con otro nombre (p. ej. eval-cinematica.tex) en
%    esta misma carpeta (el estilo y el logo se toman de ../estilo/).
% 2. Cambia los datos de la sección "DATOS DE LA EVALUACIÓN".
% 3. Compila con:  latexmk -pdf eval-cinematica.tex
%    Clave de respuestas (sin editar nada):
%    pdflatex -jobname=eval-cinematica-clave "\def\clave{}\input{eval-cinematica}"
% =====================================================================
\documentclass[11pt]{article}
% El estilo y el logo viven en la carpeta compartida plantillas/estilo/
\makeatletter\def\input@path{{../estilo/}}\makeatother
\usepackage{mmcantillo}

% ======================= DATOS DE LA EVALUACIÓN =======================
\tipodocumento{Evaluación}
\asignatura{Física}           % Cálculo, Física, Trigonometría, ...
\titulo{Cinemática en una dimensión}
\grado{10°}
\periodo{III}
\tiempo{50 min}
% \anio{2026}                 % por defecto: el año actual
% \mostrarsoluciones          % descomenta para imprimir la clave
% ======================================================================

\begin{document}

\encabezado

\begin{instrucciones}
Lee con atención cada enunciado. En la selección múltiple marca una sola
opción. En los problemas de desarrollo escribe las ecuaciones que usas y
las unidades de cada resultado. Toma $g = \qty{9,8}{m/s^2}$.
\end{instrucciones}

% ---------------------------------------------------------------------
\seccion{Selección múltiple}
% ---------------------------------------------------------------------
\begin{preguntas}

\pregunta[1] Un automóvil pasa de $\qty{0}{m/s}$ a $\qty{20}{m/s}$ en
$\qty{5}{s}$. Su aceleración media es:
\begin{opciones*}(4)
  \item $\qty{2}{m/s^2}$
  \item $\qty{4}{m/s^2}$
  \item $\qty{25}{m/s^2}$
  \item $\qty{100}{m/s^2}$
\end{opciones*}
\begin{solucion}
b) $a = \Delta v / \Delta t = 20/5 = \qty{4}{m/s^2}$.
\end{solucion}

\pregunta[1] ¿Cuál de las siguientes afirmaciones sobre la caída libre es
verdadera (sin resistencia del aire)?
\begin{opciones}
  \item Los cuerpos más pesados caen con mayor aceleración.
  \item Todos los cuerpos caen con la misma aceleración.
  \item La velocidad del cuerpo permanece constante.
  \item La aceleración es cero en el punto más alto.
\end{opciones}
\begin{solucion}
b) La aceleración es $g$ para todos los cuerpos.
\end{solucion}

\end{preguntas}

% ---------------------------------------------------------------------
\seccion{Problemas de desarrollo}
% ---------------------------------------------------------------------
\begin{preguntas}

\pregunta[3] La gráfica muestra la velocidad de un ciclista en función
del tiempo. Calcula la distancia total recorrida entre $t=0$ y
$t=\qty{10}{s}$ (el área bajo la curva).

\begin{center}
\begin{tikzpicture}
\begin{axis}[mmc, xmin=0, xmax=11, ymin=0, ymax=7,
             xlabel={$t$ (\unit{s})}, ylabel={$v$ (\unit{m/s})},
             xtick={0,2,...,10}, ytick={0,2,4,6}]
  \addplot+[name path=v, sharp plot] coordinates {(0,0) (4,6) (8,6) (10,0)};
  \path[name path=eje] (axis cs:0,0) -- (axis cs:10,0);
  \addplot[pattern=north east lines, pattern color=black!45]
    fill between[of=v and eje];
\end{axis}
\end{tikzpicture}
\end{center}
\begin{solucion}[4cm]
$d = \tfrac12(4)(6) + (4)(6) + \tfrac12(2)(6) = 12 + 24 + 6 = \qty{42}{m}$.
\end{solucion}

\pregunta[3] Se lanza una pelota verticalmente hacia arriba con velocidad
inicial $v_0 = \qty{14,7}{m/s}$.
\begin{enumerate}[label=\alph*)]
  \item ¿Cuánto tiempo tarda en alcanzar la altura máxima?
  \item ¿Cuál es la altura máxima?
\end{enumerate}
\begin{solucion}[5cm]
a) $t = v_0/g = 14{,}7/9{,}8 = \qty{1,5}{s}$.\quad
b) $h = \dfrac{v_0^2}{2g} = \dfrac{(14{,}7)^2}{2(9{,}8)} \approx \qty{11,0}{m}$.
\end{solucion}

\pregunta[2] Demuestra, derivando $x(t) = x_0 + v_0 t + \tfrac12 a t^2$,
que la aceleración es constante.
\begin{solucion}[3cm]
$\dv{x}{t} = v_0 + a t$, \quad $\dv{^2x}{t^2} = a$, que no depende de $t$.
\end{solucion}

\end{preguntas}

\vfill
\begin{center}
  \small\itshape\color{mmcGris}— ¡Éxitos! —
\end{center}

\end{document}
